\documentclass[12pt,a4paper]{article}
\usepackage[utf8]{inputenc}
\usepackage[spanish,es-tabla]{babel}
\usepackage{amsmath}
\usepackage{amsfonts}
\usepackage{amssymb}
\usepackage{graphicx}
\usepackage{tikz}
\usepackage{algorithm}
\usepackage{algpseudocode}
\usepackage{geometry}
\usepackage{fancyhdr}
\usepackage{xcolor}
\usepackage{enumitem}

\geometry{margin=2.5cm}
\pagestyle{fancy}
\fancyhf{}
\rhead{Resumen Teórico - IED}
\lhead{UCASAL}
\cfoot{\thepage}

\title{\textbf{Resumen Teórico\\Introducción a la Estructura de Datos}}
\author{Material de Estudio para Trabajo Final Integrador}
\date{Mayo 2026}

\begin{document}

\maketitle
\newpage

\section{Unidad 1: Modelos Matemáticos y Teoría de Grafos}

\subsection{Conceptos Básicos de Grafos}

\subsubsection{Definición}
Un \textbf{grafo} $G = (V, E)$ está compuesto por:
\begin{itemize}
    \item $V$: Conjunto finito no vacío de \textbf{vértices} o nodos
    \item $E$: Conjunto de \textbf{aristas} o arcos que conectan pares de vértices
\end{itemize}

\subsubsection{Tipos de Grafos}

\paragraph{Grafo Dirigido (Dígrafo):}
Las aristas tienen dirección. Se representa como $(v_i, v_j)$ donde la arista va desde $v_i$ hacia $v_j$.

\paragraph{Grafo No Dirigido:}
Las aristas no tienen dirección. La arista $\{v_i, v_j\}$ conecta ambos vértices bidireccionalmente.

\subsubsection{Propiedades de Grafos}

\paragraph{Conexidad:}
\begin{itemize}
    \item \textbf{Grafo Conexo:} Existe un camino entre cualquier par de vértices.
    \item \textbf{Grafo No Conexo:} Existen vértices sin camino entre ellos.
    \item \textbf{Componente Conexa:} Subgrafo conexo maximal.
\end{itemize}

\paragraph{Grafo Ponderado:}
Un grafo donde cada arista tiene asociado un peso o costo numérico.
$$w: E \rightarrow \mathbb{R}$$

\paragraph{Grafo Maximal y Minimal:}
\begin{itemize}
    \item \textbf{Maximal:} No se pueden agregar más aristas sin violar alguna propiedad.
    \item \textbf{Minimal:} No se pueden quitar aristas sin violar alguna propiedad.
\end{itemize}

\subsubsection{Representación de Grafos}

\paragraph{Matriz de Adyacencia:}
Matriz $A$ de tamaño $n \times n$ donde:
$$A[i][j] = \begin{cases} 
1 & \text{si existe arista de } v_i \text{ a } v_j \text{ (no ponderado)} \\
w_{ij} & \text{peso de la arista (ponderado)} \\
0 \text{ o } \infty & \text{si no existe arista}
\end{cases}$$

\textbf{Ejemplo:} Para un grafo con 4 vértices $\{A, B, C, D\}$ con aristas $A-B$, $B-C$, $C-D$:

$$\begin{bmatrix}
0 & 1 & 0 & 0 \\
1 & 0 & 1 & 0 \\
0 & 1 & 0 & 1 \\
0 & 0 & 1 & 0
\end{bmatrix}$$

\paragraph{Lista de Adyacencia:}
Para cada vértice, se mantiene una lista de sus vértices adyacentes.

\textbf{Ejemplo:}
\begin{verbatim}
A -> [B]
B -> [A, C]
C -> [B, D]
D -> [C]
\end{verbatim}

\subsection{Conceptos Adicionales}

\paragraph{Grado de un Vértice:}
\begin{itemize}
    \item \textbf{Grafo no dirigido:} Número de aristas incidentes.
    \item \textbf{Grafo dirigido:} 
    \begin{itemize}
        \item \textit{Grado de entrada:} Aristas que llegan al vértice
        \item \textit{Grado de salida:} Aristas que salen del vértice
    \end{itemize}
\end{itemize}

\paragraph{Camino:} Secuencia de vértices $v_1, v_2, ..., v_k$ donde existe arista entre vértices consecutivos.

\paragraph{Ciclo:} Camino donde el vértice inicial coincide con el final.

\newpage

\section{Unidad 2: Estructuras Lineales de Datos}

\subsection{Listas}

\subsubsection{Definición}
Una \textbf{lista} es una colección ordenada de elementos del mismo tipo, donde cada elemento tiene una posición específica.

\subsubsection{Tipos de Listas}

\paragraph{Lista Secuencial:}
Los elementos se almacenan en posiciones contiguas de memoria (array).

\textbf{Ventajas:}
\begin{itemize}
    \item Acceso directo por índice: $O(1)$
    \item Eficiente en uso de memoria
\end{itemize}

\textbf{Desventajas:}
\begin{itemize}
    \item Inserción/eliminación costosa: $O(n)$
    \item Tamaño fijo
\end{itemize}

\paragraph{Lista Enlazada (Linkeada):}
Los elementos se almacenan en nodos, cada uno con un puntero al siguiente.

\textbf{Estructura de un Nodo:}
\begin{verbatim}
struct Nodo {
    dato: TipoDato
    siguiente: Puntero a Nodo
}
\end{verbatim}

\textbf{Ventajas:}
\begin{itemize}
    \item Inserción/eliminación eficiente: $O(1)$ si se tiene el puntero
    \item Tamaño dinámico
\end{itemize}

\textbf{Desventajas:}
\begin{itemize}
    \item Acceso secuencial: $O(n)$
    \item Mayor uso de memoria (punteros)
\end{itemize}

\subsection{Pila (Stack)}

\subsubsection{Definición}
Estructura LIFO (\textit{Last In, First Out}): el último elemento en entrar es el primero en salir.

\subsubsection{Operaciones Básicas}
\begin{itemize}
    \item \texttt{Push(x)}: Agregar elemento $x$ al tope
    \item \texttt{Pop()}: Eliminar y retornar el elemento del tope
    \item \texttt{Top()} o \texttt{Peek()}: Consultar el elemento del tope sin eliminarlo
    \item \texttt{isEmpty()}: Verificar si la pila está vacía
\end{itemize}

\subsubsection{Ejemplo Práctico}
\begin{verbatim}
Operaciones:        Pila resultante:
Push(5)             [5]
Push(3)             [5, 3]
Push(8)             [5, 3, 8]  <- Tope
Pop()               [5, 3]     (retorna 8)
Push(1)             [5, 3, 1]  <- Tope
\end{verbatim}

\subsection{Cola (Queue)}

\subsubsection{Definición}
Estructura FIFO (\textit{First In, First Out}): el primer elemento en entrar es el primero en salir.

\subsubsection{Operaciones Básicas}
\begin{itemize}
    \item \texttt{Enqueue(x)} o \texttt{Insert(x)}: Agregar elemento $x$ al final
    \item \texttt{Dequeue()} o \texttt{Remove()}: Eliminar y retornar el elemento del frente
    \item \texttt{Front()}: Consultar el elemento del frente sin eliminarlo
    \item \texttt{isEmpty()}: Verificar si la cola está vacía
\end{itemize}

\subsubsection{Ejemplo Práctico}
\begin{verbatim}
Operaciones:           Cola resultante:
Enqueue(7)             [7]
Enqueue(3)             [7, 3]
Enqueue(9)             [7, 3, 9]
Dequeue()              [3, 9]      (retorna 7)
Enqueue(2)             [3, 9, 2]
                       ^Frente    ^Final
\end{verbatim}

\subsection{Asignación Dinámica de Memoria}

\paragraph{Concepto:}
Permite solicitar memoria durante la ejecución del programa según sea necesario.

\paragraph{Operaciones:}
\begin{itemize}
    \item \texttt{malloc()} / \texttt{new}: Reservar memoria
    \item \texttt{free()} / \texttt{delete}: Liberar memoria
\end{itemize}

\subsection{Lista Irrestricta}

\paragraph{Definición:}
Lista donde se pueden realizar operaciones de inserción y eliminación en cualquier posición, no solo en los extremos.

\paragraph{Operaciones:}
\begin{itemize}
    \item Insertar en posición $i$
    \item Eliminar en posición $i$
    \item Buscar elemento
    \item Recorrer la lista
\end{itemize}

\newpage

\section{Unidad 3: Árboles}

\subsection{Definiciones Básicas}

\subsubsection{Árbol}
Estructura jerárquica compuesta por nodos conectados por aristas, donde:
\begin{itemize}
    \item Existe un nodo especial llamado \textbf{raíz}
    \item Cada nodo puede tener cero o más \textbf{hijos}
    \item No existen ciclos
\end{itemize}

\subsubsection{Terminología}
\begin{itemize}
    \item \textbf{Raíz:} Nodo superior sin padre
    \item \textbf{Hoja:} Nodo sin hijos
    \item \textbf{Nodo Interno:} Nodo con al menos un hijo
    \item \textbf{Padre:} Nodo predecesor directo
    \item \textbf{Hijo:} Nodo sucesor directo
    \item \textbf{Hermanos:} Nodos con el mismo padre
    \item \textbf{Subárbol:} Árbol formado por un nodo y todos sus descendientes
    \item \textbf{Nivel:} Distancia desde la raíz (la raíz está en nivel 0)
    \item \textbf{Altura:} Máximo nivel del árbol (longitud del camino más largo desde la raíz a una hoja)
    \item \textbf{Profundidad de un nodo:} Longitud del camino desde la raíz hasta ese nodo
\end{itemize}

\subsection{Árboles Binarios}

\subsubsection{Definición}
Árbol donde cada nodo tiene \textbf{como máximo dos hijos}: hijo izquierdo e hijo derecho.

\subsubsection{Tipos de Árboles Binarios}

\paragraph{Árbol Binario Completo:}
Todos los niveles están completamente llenos, excepto posiblemente el último, que se llena de izquierda a derecha.

\paragraph{Árbol Binario Lleno:}
Todos los nodos internos tienen exactamente 2 hijos.

\paragraph{Árbol Binario Perfecto:}
Todos los niveles están completamente llenos. Tiene $2^{h+1} - 1$ nodos, donde $h$ es la altura.

\paragraph{Árbol Binario Degenerado:}
Cada nodo tiene solo un hijo (similar a una lista enlazada).

\subsubsection{Propiedades}
\begin{itemize}
    \item Número máximo de nodos en nivel $i$: $2^i$
    \item Número máximo de nodos en árbol de altura $h$: $2^{h+1} - 1$
    \item Altura mínima con $n$ nodos: $\lceil \log_2(n+1) \rceil - 1$
\end{itemize}

\subsection{Recorridos de Árboles Binarios}

\subsubsection{Pre-orden (Raíz-Izquierda-Derecha)}
\begin{algorithm}[H]
\caption{Recorrido Pre-orden}
\begin{algorithmic}
\Procedure{PreOrden}{nodo}
    \If{nodo $\neq$ null}
        \State Visitar(nodo)
        \State PreOrden(nodo.izquierdo)
        \State PreOrden(nodo.derecho)
    \EndIf
\EndProcedure
\end{algorithmic}
\end{algorithm}

\subsubsection{In-orden (Izquierda-Raíz-Derecha)}
\begin{algorithm}[H]
\caption{Recorrido In-orden}
\begin{algorithmic}
\Procedure{InOrden}{nodo}
    \If{nodo $\neq$ null}
        \State InOrden(nodo.izquierdo)
        \State Visitar(nodo)
        \State InOrden(nodo.derecho)
    \EndIf
\EndProcedure
\end{algorithmic}
\end{algorithm}

\subsubsection{Post-orden (Izquierda-Derecha-Raíz)}
\begin{algorithm}[H]
\caption{Recorrido Post-orden}
\begin{algorithmic}
\Procedure{PostOrden}{nodo}
    \If{nodo $\neq$ null}
        \State PostOrden(nodo.izquierdo)
        \State PostOrden(nodo.derecho)
        \State Visitar(nodo)
    \EndIf
\EndProcedure
\end{algorithmic}
\end{algorithm}

\subsubsection{Ejemplo de Recorridos}
Para el árbol:
\begin{verbatim}
      A
     / \
    B   C
   / \
  D   E
\end{verbatim}

\begin{itemize}
    \item \textbf{Pre-orden:} A, B, D, E, C
    \item \textbf{In-orden:} D, B, E, A, C
    \item \textbf{Post-orden:} D, E, B, C, A
\end{itemize}

\subsection{Árboles Balanceados}

\subsubsection{Definición}
Un árbol está \textbf{balanceado} si para cada nodo, la diferencia de altura entre su subárbol izquierdo y derecho es como máximo 1.

\subsubsection{Factor de Balance}
$$FB(nodo) = \text{altura(subárbol\_izquierdo)} - \text{altura(subárbol\_derecho)}$$

Un árbol está balanceado si: $|FB(nodo)| \leq 1$ para todo nodo.

\subsection{Árboles N-arios}

\subsubsection{Definición}
Árbol donde cada nodo puede tener un número arbitrario de hijos.

\subsubsection{Transformada de Knuth}
Técnica para convertir un árbol N-ario en un árbol binario:
\begin{itemize}
    \item El \textbf{primer hijo} de un nodo se convierte en su \textbf{hijo izquierdo}
    \item El \textbf{siguiente hermano} se convierte en su \textbf{hijo derecho}
\end{itemize}

\textbf{Ejemplo:}
\begin{verbatim}
Árbol N-ario:        Árbol Binario (Knuth):
      A                    A
    / | \                 /
   B  C  D               B
  /|                      \
 E F                       C
                            \
                             D
                            /
                           E
                            \
                             F
\end{verbatim}

\newpage

\section{Unidad 4: Acceso de Datos}

\subsection{Tablas de Símbolos}

\subsubsection{Definición}
Estructura de datos que permite almacenar pares \textit{(clave, valor)} con las operaciones:
\begin{itemize}
    \item \texttt{Insert(clave, valor)}
    \item \texttt{Search(clave)} $\rightarrow$ valor
    \item \texttt{Delete(clave)}
\end{itemize}

\subsubsection{Costos de Acceso}
\begin{itemize}
    \item \textbf{Búsqueda secuencial:} $O(n)$
    \item \textbf{Búsqueda binaria (arreglo ordenado):} $O(\log n)$
    \item \textbf{Árbol de búsqueda:} $O(h)$ donde $h$ es la altura
    \item \textbf{Tabla Hash:} $O(1)$ promedio
\end{itemize}

\subsection{Árboles de Búsqueda Binaria (ABB)}

\subsubsection{Definición}
Árbol binario donde para cada nodo:
\begin{itemize}
    \item Todos los valores del subárbol izquierdo son \textbf{menores}
    \item Todos los valores del subárbol derecho son \textbf{mayores}
\end{itemize}

\subsubsection{Operaciones}

\paragraph{Búsqueda:}
\begin{algorithm}[H]
\caption{Búsqueda en ABB}
\begin{algorithmic}
\Procedure{Buscar}{nodo, clave}
    \If{nodo = null}
        \State \Return null
    \EndIf
    \If{clave = nodo.valor}
        \State \Return nodo
    \ElsIf{clave < nodo.valor}
        \State \Return Buscar(nodo.izquierdo, clave)
    \Else
        \State \Return Buscar(nodo.derecho, clave)
    \EndIf
\EndProcedure
\end{algorithmic}
\end{algorithm}

\paragraph{Inserción:}
Similar a búsqueda, pero se agrega el nodo en la posición correspondiente.

\paragraph{Eliminación:}
Casos:
\begin{enumerate}
    \item Nodo hoja: Se elimina directamente
    \item Nodo con un hijo: Se reemplaza por su hijo
    \item Nodo con dos hijos: Se reemplaza por su predecesor o sucesor in-orden
\end{enumerate}

\subsection{Árboles AVL}

\subsubsection{Definición}
ABB autobalanceado donde el factor de balance de cada nodo es $-1$, $0$ o $1$.

\subsubsection{Rotaciones}
Para mantener el balance después de inserciones/eliminaciones:

\paragraph{Rotación Simple a la Derecha:}
Usada cuando el subárbol izquierdo del hijo izquierdo causa desbalance.

\paragraph{Rotación Simple a la Izquierda:}
Usada cuando el subárbol derecho del hijo derecho causa desbalance.

\paragraph{Rotación Doble Izquierda-Derecha:}
Usada cuando el subárbol derecho del hijo izquierdo causa desbalance.

\paragraph{Rotación Doble Derecha-Izquierda:}
Usada cuando el subárbol izquierdo del hijo derecho causa desbalance.

\subsection{Hashing}

\subsubsection{Concepto}
Técnica que utiliza una función hash para mapear claves a posiciones en una tabla.

$$h: K \rightarrow \{0, 1, ..., m-1\}$$

donde $K$ es el conjunto de claves y $m$ es el tamaño de la tabla.

\subsubsection{Funciones de Transformación}

\paragraph{Método de División:}
$$h(k) = k \bmod m$$

donde $m$ es generalmente un número primo.

\paragraph{Método de Multiplicación:}
$$h(k) = \lfloor m \cdot (k \cdot A \bmod 1) \rfloor$$

donde $0 < A < 1$ (comúnmente $A \approx 0.618$).

\paragraph{Para Cadenas de Caracteres:}
Convertir cada carácter a su código ASCII y aplicar una función:
$$h(s) = \left(\sum_{i=0}^{n-1} ASCII(s[i]) \cdot 31^i\right) \bmod m$$

\textbf{Ejemplo con método de división:}
\begin{verbatim}
Clave: "Hola"
Conversión: H(72) + o(111) + l(108) + a(97) = 388
Con m = 50: h("Hola") = 388 mod 50 = 38
\end{verbatim}

\subsubsection{Colisiones}

\paragraph{Definición:}
Ocurre cuando dos claves diferentes producen el mismo valor hash: $h(k_1) = h(k_2)$ con $k_1 \neq k_2$.

\paragraph{Resolución por Encadenamiento:}
Cada posición de la tabla contiene una lista enlazada con todos los elementos que tienen ese hash.

\paragraph{Resolución por Direccionamiento Abierto:}
\begin{itemize}
    \item \textbf{Prueba Lineal:} $h'(k, i) = (h(k) + i) \bmod m$
    \item \textbf{Prueba Cuadrática:} $h'(k, i) = (h(k) + c_1 \cdot i + c_2 \cdot i^2) \bmod m$
    \item \textbf{Hash Doble:} $h'(k, i) = (h_1(k) + i \cdot h_2(k)) \bmod m$
\end{itemize}

\subsection{Árboles B}

\subsubsection{Definición}
Árbol de búsqueda N-ario balanceado optimizado para acceso a disco, donde:
\begin{itemize}
    \item Cada nodo puede contener múltiples claves
    \item Todas las hojas están al mismo nivel
    \item Se mantiene un orden de $t$ (grado mínimo)
\end{itemize}

\subsubsection{Propiedades del Árbol B de orden $t$:}
\begin{itemize}
    \item Cada nodo (excepto la raíz) tiene entre $t-1$ y $2t-1$ claves
    \item Cada nodo interno (excepto la raíz) tiene entre $t$ y $2t$ hijos
    \item La raíz tiene al menos 1 clave
    \item Todas las hojas están al mismo nivel
\end{itemize}

\newpage

\section{Ejemplos Prácticos Adicionales}

\subsection{Ejemplo: Operaciones con Pila y Cola}

\textbf{Situación:} Dada la lista: [12, 7, 18, 5, 9, 14]

\textbf{Proceso:}
\begin{itemize}
    \item Elementos pares (12, 18, 14) $\rightarrow$ Pila: [12, 18, 14] (tope en 14)
    \item Elementos impares (7, 5, 9) $\rightarrow$ Cola: [7, 5, 9]
    \item Promedio de la pila: $(12 + 18 + 14) / 3 = 14.67$
\end{itemize}

\subsection{Ejemplo: Recorrido de Árbol}

Para el árbol:
\begin{verbatim}
        15
       /  \
      10   20
     / \   / \
    5  12 18 25
\end{verbatim}

\begin{itemize}
    \item \textbf{Tipo:} Árbol binario de búsqueda
    \item \textbf{Completo:} Sí
    \item \textbf{Altura:} 2
    \item \textbf{Balanceado:} Sí (FB = 0 para todos los nodos)
    \item \textbf{Pre-orden:} 15, 10, 5, 12, 20, 18, 25
    \item \textbf{In-orden:} 5, 10, 12, 15, 18, 20, 25
    \item \textbf{Post-orden:} 5, 12, 10, 18, 25, 20, 15
\end{itemize}

\subsection{Ejemplo: Hashing con Método de División}

\textbf{Claves:} ``Ana'', ``Juan'', ``María'', ``Pedro''\\
\textbf{Tabla:} $m = 50$

\textbf{Conversión ASCII y cálculo:}
\begin{align*}
\text{Ana} &= 65 + 110 + 97 = 272 \rightarrow 272 \bmod 50 = 22\\
\text{Juan} &= 74 + 117 + 97 + 110 = 398 \rightarrow 398 \bmod 50 = 48\\
\text{María} &= 77 + 97 + 114 + 237 + 97 = 622 \rightarrow 622 \bmod 50 = 22\\
\text{Pedro} &= 80 + 101 + 100 + 114 + 111 = 506 \rightarrow 506 \bmod 50 = 6
\end{align*}

\textbf{Nota:} ``Ana'' y ``María'' tienen colisión en la posición 22.

\newpage

\section{Algoritmos Fundamentales}

\subsection{Algoritmo de Recorrido de Lista con Pila y Cola}

\begin{algorithm}[H]
\caption{Procesar Lista en Pila y Cola}
\begin{algorithmic}
\Procedure{ProcesarLista}{lista}
    \State pila $\leftarrow$ nueva Pila()
    \State cola $\leftarrow$ nueva Cola()
    \State suma\_pila $\leftarrow$ 0
    \State contador\_pila $\leftarrow$ 0
    
    \For{cada elemento en lista}
        \If{elemento mod 2 = 0} \Comment{Es par}
            \State pila.push(elemento)
            \State suma\_pila $\leftarrow$ suma\_pila + elemento
            \State contador\_pila $\leftarrow$ contador\_pila + 1
        \Else \Comment{Es impar}
            \State cola.enqueue(elemento)
        \EndIf
    \EndFor
    
    \State Mostrar(pila)
    \State Mostrar(cola)
    
    \If{contador\_pila > 0}
        \State promedio $\leftarrow$ suma\_pila / contador\_pila
        \State Mostrar(``Promedio de pila:'', promedio)
    \EndIf
\EndProcedure
\end{algorithmic}
\end{algorithm}

\subsection{Verificación de Balance en Árbol}

\begin{algorithm}[H]
\caption{Verificar si un árbol está balanceado}
\begin{algorithmic}
\Procedure{CalcularAltura}{nodo}
    \If{nodo = null}
        \State \Return -1
    \EndIf
    \State \Return $1 + \max($CalcularAltura(nodo.izq), CalcularAltura(nodo.der)$)$
\EndProcedure

\Procedure{EstaBalanceado}{nodo}
    \If{nodo = null}
        \State \Return true
    \EndIf
    
    \State altura\_izq $\leftarrow$ CalcularAltura(nodo.izquierdo)
    \State altura\_der $\leftarrow$ CalcularAltura(nodo.derecho)
    \State factor\_balance $\leftarrow$ altura\_izq - altura\_der
    
    \If{$|$factor\_balance$| > 1$}
        \State \Return false
    \EndIf
    
    \State \Return EstaBalanceado(nodo.izq) AND EstaBalanceado(nodo.der)
\EndProcedure
\end{algorithmic}
\end{algorithm}

\newpage

\section{Complejidades Temporales - Resumen}

\begin{center}
\begin{tabular}{|l|c|c|c|}
\hline
\textbf{Estructura} & \textbf{Búsqueda} & \textbf{Inserción} & \textbf{Eliminación} \\
\hline
Array (desordenado) & $O(n)$ & $O(1)$ & $O(n)$ \\
Array (ordenado) & $O(\log n)$ & $O(n)$ & $O(n)$ \\
Lista Enlazada & $O(n)$ & $O(1)^*$ & $O(1)^*$ \\
Pila & - & $O(1)$ & $O(1)$ \\
Cola & - & $O(1)$ & $O(1)$ \\
ABB (balanceado) & $O(\log n)$ & $O(\log n)$ & $O(\log n)$ \\
ABB (degenerado) & $O(n)$ & $O(n)$ & $O(n)$ \\
AVL & $O(\log n)$ & $O(\log n)$ & $O(\log n)$ \\
Hash (ideal) & $O(1)$ & $O(1)$ & $O(1)$ \\
Árbol B & $O(\log n)$ & $O(\log n)$ & $O(\log n)$ \\
\hline
\end{tabular}
\end{center}

\small{*Si se tiene el puntero al nodo}

\section{Consejos para el Trabajo Integrador}

\subsection{Análisis de Grafos}
\begin{enumerate}
    \item Identificar vértices y aristas
    \item Verificar dirección de las aristas
    \item Comprobar conectividad (¿hay camino entre todos los pares?)
    \item Verificar si hay pesos en las aristas
    \item Construir la matriz siguiendo el orden de los vértices
\end{enumerate}

\subsection{Trabajar con Pilas y Colas}
\begin{enumerate}
    \item Identificar claramente el criterio de separación
    \item Mantener el orden LIFO para pilas (último que entra, primero que sale)
    \item Mantener el orden FIFO para colas (primero que entra, primero que sale)
    \item Para promedios: sumar todos los elementos y dividir por la cantidad
\end{enumerate}

\subsection{Análisis de Árboles}
\begin{enumerate}
    \item Identificar la raíz y los nodos hoja
    \item Calcular la altura contando niveles desde 0
    \item Para verificar balance: calcular altura de subárboles izquierdo y derecho
    \item Para recorridos: seguir estrictamente el orden indicado (pre/in/post)
    \item Un árbol es completo si todos los niveles están llenos excepto posiblemente el último
\end{enumerate}

\subsection{Hashing}
\begin{enumerate}
    \item Convertir cada carácter a su valor ASCII
    \item Sumar todos los valores
    \item Aplicar módulo $m$ (el tamaño de la tabla)
    \item Documentar claramente los cálculos
    \item Identificar posibles colisiones (mismo resultado para claves diferentes)
\end{enumerate}

\vfill

\begin{center}
\textit{Este resumen cubre todos los conceptos teóricos necesarios\\para resolver el Trabajo Final Integrador de Introducción a la Estructura de Datos.}

\vspace{1em}

\textbf{¡Éxitos en el examen!}
\end{center}

\end{document}